\section{Visión general del sistema de conocimientos del curso}

En la conclusión del curso, el ponente repasa la evolución del conocimiento acumulado hasta el momento.

\subsection{Trayectoria de evolución tecnológica}

Desde el inicio del curso hasta la décima clase, el ponente ha construido la siguiente trayectoria tecnológica completa:

\begin{enumerate}
    \item \textbf{VAE/GAN} (Lecture 1--2): modelos generativos de un solo paso, limitados en calidad y diversidad.
    \item \textbf{DDPM} (Lecture 3--4): introducción del desruido iterativo multietapa, lo que eleva drásticamente la calidad generativa.
    \item \textbf{DDIM} (Lecture 5--6): variante de muestreo determinista de DDPM, compatible con saltos en pasos.
    \item \textbf{SDE/ODE} (Lecture 7): marco de tiempo continuo, perspectiva unificada.
    \item \textbf{DPM-Solver} (Lecture 8): aceleración del muestreo aprovechando la estructura ODE.
    \item \textbf{Consistency Model} (Lecture 9): aprendizaje de mapas rápidos para lograr generación con pocos pasos.
    \item \textbf{Flow Matching} (Lecture 9--10): marco más general; Rectified Flow + Reflow logra trayectorias rectas y generación en un paso.
\end{enumerate}

\begin{warningbox}{Flow Matching no es una ``alternativa'' a los modelos de difusión}
Aunque esta clase destaca las ventajas de los modelos de flujo, es necesario observar lo siguiente:
\begin{itemize}
    \item Los modelos de difusión y los modelos de flujo son matemáticamente equivalentes (bajo trayectorias condicionales gaussianas lineales).
    \item Un modelo de difusión ya entrenado puede convertirse directamente en un modelo de flujo para su uso.
    \item La elección del marco depende principalmente de consideraciones prácticas: estabilidad del entrenamiento, eficiencia del muestreo y complejidad del código.
    \item Las ``mejores prácticas'' actuales consisten en utilizar el marco Flow Matching + entrenamiento con Rectified Flow + ajuste fino (fine-tuning) con Reflow.
\end{itemize}
\end{warningbox}

\subsection{Próximamente: contenido de las siguientes clases}

El ponente anuncia el contenido de las tres próximas clases:
\begin{itemize}
    \item \textbf{Próxima semana}: Inference-time Guidance (guía en tiempo de inferencia), cómo inyectar preferencias del usuario sin volver a entrenar.
    \item \textbf{Posteriormente}: Discrete Diffusion Models (modelos de difusión para datos discretos), extensión de las ideas de difusión/flujo a datos no continuos (por ejemplo, texto).
\end{itemize}

\subsection{Resumen de la sección}

Desde GAN hasta los modelos de difusión y luego los modelos de flujo, el campo de los modelos generativos ha experimentado una evolución desde un paso, pasando por múltiples pasos, hasta pocos pasos. Flow Matching ofrece un marco unificado y más flexible que permite comprender los logros existentes de los modelos de difusión al mismo tiempo que abre espacio para mejoras futuras.
